We fix conventions: graphs are finite, simple, undirected.
For a graph $G$, let $\Delta(G)$ denote its maximum degree.
For two edges $e,f\in E(G)$ define their \emph{edge-distance}
\[
\operatorname{dist}_G(e,f):=\min\{\operatorname{dist}_G(u,v): u\in e,\ v\in f\},
\]
where $\operatorname{dist}_G(u,v)$ is the usual shortest-path distance between vertices.
Thus $\operatorname{dist}_G(e,f)=0$ iff $e$ and $f$ share a vertex.

For integers $t\ge 1$ and $d\ge 1$, define $h_t(d)$ to be the least integer $m$ such that every graph $G$ with $\Delta(G)\le d$ and $|E(G)|\ge m$ contains two edges $e,f$ with $\operatorname{dist}_G(e,f)\ge t$.
Equivalently, $h_t(d)-1$ is the maximum number of edges in a graph with maximum degree $\le d$ and with \emph{edge-diameter} at most $t-1$.
